% !TeX root = ../main.tex
% ---------------------------------------------------------------------------
% Groundwater and Hydrogeology for Underground Works
% Scope extension: not derived from the RSE3010 weekly material.
% Key sources: Goodman et al. (1965); Heuer (1995, 2005); El Tani (2003);
% Freeze & Cherry; ITA WG on drainage and waterproofing.
% ---------------------------------------------------------------------------
\chapter{Groundwater and Hydrogeology for Underground Works}
\label{ch:hydrogeology}


\section{Groundwater in underground design}
\label{sec:groundwater-in-underground-design}

\subsection{Why groundwater governs underground works}
\label{sec:groundwater-governs}
\scopenote{Water as the single most common cause of cost and programme
overrun. Its four distinct roles: effective stress, inflow, face instability,
and long-term loading on the lining.}

\subsection{Conceptual hydrogeological models}
\label{sec:hydrogeological-model}
\scopenote{Aquifers, aquitards, perched water, recharge and boundary
conditions. Fracture-controlled versus porous-medium flow, and when the
equivalent-continuum assumption breaks down.}

\section{Permeability and pore pressure}
\label{sec:permeability-and-pore-pressure}

\subsection{Permeability and its measurement}
\label{sec:permeability}
\scopenote{Laboratory permeability, falling head, packer and Lugeon tests,
pumping tests and their differing support volumes. The scale dependence of
measured $k$, and why single-borehole values mislead.}

\subsection{Effective stress and pore pressure around an opening}
\label{sec:effective-stress-tunnel}
\scopenote{Pore pressure redistribution on excavation; drained and undrained
response; the effective stress path at the face. Links to
Chapter~\ref{ch:soft-ground}.}

\section{Inflow prediction}
\label{sec:inflow-prediction}
\scopenote{Analytical steady-state solutions (Goodman, El Tani), Heuer's
correction for fracture flow, probabilistic inflow along an alignment, and
numerical seepage analysis. Realistic accuracy expectations.}

\section{Drainage philosophy and lining loads}
\label{sec:drainage-philosophy-and-lining-loads}

\subsection{Drained versus tanked design philosophy}
\label{sec:drained-vs-tanked}
\scopenote{The fundamental lining decision: relieve water pressure and accept
permanent drainage and maintenance, or resist full hydrostatic head. Cost,
durability, environmental and operational consequences of each.}

\subsection{Water pressure and uplift on linings}
\label{sec:water-pressure-lining}
\scopenote{Design head selection, long-term recovery of the water table,
buoyancy of shallow tunnels and shafts, and load cases for
Chapters~\ref{ch:segmental-linings} and~\ref{ch:final-linings}.}

\section{Environmental effects and groundwater control}
\label{sec:environmental-effects-and-groundwater-control}

\subsection{Drawdown, consolidation settlement and environmental impact}
\label{sec:drawdown}
\scopenote{Settlement caused by dewatering rather than by volume loss; impact
on groundwater-dependent ecosystems, bores and existing foundations; the
regulatory framing under Victorian and Commonwealth approvals.}

\subsection{Groundwater control methods}
\label{sec:groundwater-control}
\scopenote{Dewatering, cut-off walls, pre-excavation grouting, ground freezing
and compressed air, with selection criteria by permeability and head. Detail
in Chapter~\ref{ch:ground-improvement}.}

\section{Inrush and water chemistry}
\label{sec:inrush-and-water-chemistry}

\subsection{Water inrush and high-pressure features}
\label{sec:inrush}
\scopenote{Karst conduits, fault-hosted aquifers and confined high-head
features; probe drilling as standard practice. Hazard management in
Chapter~\ref{ch:construction-hazards}.}

\subsection{Water chemistry, acid rock drainage and discharge}
\label{sec:water-chemistry}
\scopenote{Aggressive groundwater and its effect on concrete durability;
sulfide oxidation and acid rock drainage; treatment and discharge compliance.}
